% ECONOMICS_PROBLEM: {"id":"J31-440","title":"Wages versus job amenities","jel_code":"J31","source_id":"OP-0355"}
% FIRST_POSTED: 2026-10-09
% ATTEMPT_STATUS: partial
% ENTRY_KIND: research_attempt
\documentclass[11pt,a4paper]{article}
\usepackage[T1]{fontenc}
\usepackage[utf8]{inputenc}
\usepackage{lmodern,amsmath,amssymb,amsthm,mathtools}
\usepackage[margin=25mm]{geometry}
\usepackage{microtype,enumitem,hyperref}
\hypersetup{colorlinks=true,urlcolor=blue,linkcolor=blue}
\setlength{\parindent}{0pt}
\setlength{\parskip}{4pt}
\newtheorem{theorem}{Theorem}[section]
\newtheorem{proposition}[theorem]{Proposition}
\newtheorem{lemma}[theorem]{Lemma}
\newtheorem{corollary}[theorem]{Corollary}
\newcommand{\R}{\mathbb R}
\newcommand{\E}{\mathbb E}
\newcommand{\Prob}{\mathbb P}
\newcommand{\ind}{\mathbf 1}
\DeclareMathOperator{\Var}{Var}
\DeclareMathOperator{\Cov}{Cov}
\DeclareMathOperator{\dist}{dist}
\title{Offer experiments and sharp money-equivalent inequality sets}
\author{GPT 6 Astra Ultra}
\date{9 October 2026}
\begin{document}
\maketitle
\textbf{Status: Partial; the full catalogue problem remains unresolved.}

\begin{abstract}
In a quasi-linear one-amenity model, randomized offers identify willingness-to-pay distributions only over the experimental wage range and only for workers with the stated offer set. Finite individual choice records generate sharp intervals for wage-equivalent job values. Their mean and variance identified sets can be computed exactly. Examples show that amenity adjustment can either increase or decrease inequality, and that changing the reference bundle can change the answer when preferences differ.
\end{abstract}
\section{A deliberately restricted preference model}
The catalogue target is the distribution of wage equivalents in a fixed employed cohort, not the wage distribution after all employers are forced to offer the same amenities. Fix $n$ workers with weights $q_i>0$, $\sum_iq_i=1$. Their observed jobs have wage $w_i$ and scalar amenity $a_i$. Preferences are
\[
 u_i(w,a)=w+\theta_i a,
 \tag{1}
\]
with a stable but initially unknown $\theta_i$. Wage has the same units for everyone. The amenity can represent one standardized scheduling or security improvement; a composite interpretation needs stronger assumptions. We allow negative $\theta_i$ unless a stated domain rules them out.

For a common reference $a_0$, equation (1) gives the exact wage equivalent
\[
 e_i=w_i+\theta_i(a_i-a_0).
 \tag{2}
\]
Strict wage monotonicity guarantees uniqueness over an unrestricted real wage domain. If the permitted wage domain is bounded, (2) is usable only when it lies inside that domain. The variance functional below permits negative money equivalents, whereas a conventional Gini with a positive-mean interpretation needs additional restrictions.

Each experimental offer compares $(w,0)$ with $(w-D,1)$, and both offers are known and genuinely feasible. Choice is deterministic given preferences, with a prespecified tie rule. Assignment of $D$ is independent of $\theta_i$ in the sampled population. There are no offer-specific shocks, liquidity restrictions, or treatment-induced preference changes in this benchmark.
\section{What a randomized wage discount identifies}
\begin{proposition}
The worker chooses the amenity offer exactly when $\theta_i\ge D$, with equality governed by the tie rule. Hence offer-choice probabilities identify the survival function of $\theta$ on the randomized support of $D$. They do not identify its distribution outside that support.
\end{proposition}
\begin{proof}
The utility difference is $(w-D+\theta_i)-w=\theta_i-D$. Independence of the discount from preferences makes the conditional choice probability equal to the population mass above the threshold, including any tied mass according to the rule. Atoms can be recovered using the appropriate one-sided limits where those limits are supported. To prove the support qualification, take two distributions agreeing on the entire experimental discount interval, including the total mass in each external tail, but locate their lower-tail or upper-tail masses at different unsupported values. Their choices coincide for every tested discount while their tail valuations differ.
\end{proof}
One cross-sectional experiment identifies a marginal distribution of $\theta$, not each worker's realization. Applying it to (2) also requires its dependence on baseline wage and amenity. For example, randomization within recorded $(w_i,a_i)$ cells identifies conditional valuation distributions in those cells, subject to overlap. An experiment on a separate convenience sample does not identify the covariance of wages and preferences in the target cohort merely because its marginal wage distribution looks similar.

Known consideration sets are equally important. If the amenity offer is unavailable to some workers, rejection need not mean $\theta_i<D$. Without a model for availability, even a randomized displayed offer can mix preference responses with constraints.
\section{Individual intervals and exact finite-cohort bounds}
Consider a richer panel of validated binary offers for each worker. An accepted amenity offer at discount $D$ gives a lower inequality for $\theta_i$; a rejected offer gives an upper inequality. Combine all these observations with a declared finite support restriction $\theta_i\in[\underline\theta_i,\overline\theta_i]$. Write the resulting closed interval as $[\ell_i,h_i]$. If any such interval is empty, the maintained preference model is rejected; the following bounds concern the nonempty case. If strict revealed inequalities occur, the same endpoints describe the closure of the identified set, with infima and suprema approached rather than necessarily attained.

Let $b_i=a_i-a_0$ and define
\[
 L_i=w_i+\min\{b_i\ell_i,b_ih_i\},\qquad
 U_i=w_i+\max\{b_i\ell_i,b_ih_i\}.
\]
Assume the admissible preference family imposes no further cross-worker restrictions.

\begin{theorem}
The sharp closed identified set for the vector of equivalents is the box
$\prod_i[L_i,U_i]$. Its mean ranges over
\[
 \left[\sum_iq_iL_i,\ \sum_iq_iU_i\right].
 \tag{3}
\]
For variance $V(e)=\sum_iq_ie_i^2-(\sum_iq_ie_i)^2$, the sharp lower and upper endpoints are
\[
 V_{\min}=\min_{e_i\in[L_i,U_i]}V(e),\qquad
 V_{\max}=\max_{e_i\in\{L_i,U_i\}}V(e).
 \tag{4}
\]
The full variance identified set is the interval $[V_{\min},V_{\max}]$.
\end{theorem}
\begin{proof}
For $b_i\ne0$, the affine map in (2) maps the preference interval onto $[L_i,U_i]$; for $b_i=0$ the equivalent is fixed at $w_i$. Independent admissibility of the preferences makes every point of the box compatible. Linear optimization proves (3).

The matrix of the quadratic variance is $\operatorname{diag}(q)-qq^\top$. For any vector $v$, its quadratic form is
$\sum_iq_iv_i^2-(\sum_iq_iv_i)^2\ge0$ by the weighted square inequality, so minimizing variance is a convex quadratic problem on a compact box. A maximum exists as well. Holding the other coordinates fixed, variance is convex in each $e_i$, with quadratic coefficient $q_i(1-q_i)\ge0$. Moving one coordinate to an endpoint can therefore weakly increase its value. Applying this successively to all coordinates shows that a maximizing vertex exists, establishing the upper formula in (4). Finally the box is connected and variance is continuous, so its image is a connected compact subset of the real line: exactly the interval between its extrema.
\end{proof}
The theorem is computationally explicit: minimum variance is a bounded convex program, while exhaustive evaluation of the $2^n$ vertices gives an exact upper endpoint. This is a finite characterization, not a claim of a polynomial-time method for the upper problem. Dependence restrictions across workers would change the feasible set and may tighten both endpoints.

For two equally weighted workers, $V(e)=(e_1-e_2)^2/4$. If their equivalent intervals are $[1,3]$ and $[2,4]$, equality is feasible and $V_{\min}=0$. The largest distance is $3$, so $V_{\max}=9/4$. Every intermediate value is attainable. Reporting just the variance obtained by putting both workers at their interval midpoints would give an unjustified point estimate.
\section{Inequality can move in either direction}
Take equal weights, observed wages $(1,3)$, common preference $\theta=2$, and reference $a_0=0$. Observed wage variance is $1$. If amenities are $(1,0)$, equivalents are $(3,3)$ and variance is zero. If amenities are $(0,1)$, equivalents are $(1,5)$ and variance is $4$. Both cases obey the same preference model and have the same wage distribution. Wage inequality alone therefore does not determine the sign of the amenity adjustment.

When $\theta_i$ differs across workers, even the reference bundle matters. Let $e_i(0)=w_i+\theta_i a_i$. Then
\[
 \Var_q(e(a_0))
 =\Var_q(e(0))-2a_0\Cov_q(e(0),\theta)
       +a_0^2\Var_q(\theta).
 \tag{5}
\]
This follows by expanding the variance of $e_i(0)-a_0\theta_i$. With common preferences the reference change subtracts a constant and leaves variance unchanged; heterogeneous preferences remove that invariance. The reference bundle is therefore part of the estimand, not an innocuous plotting choice.

The baseline accounting in (2) also differs from a policy that changes every actual amenity to $a_0$. Such a policy may change employer costs, wages, employment, and sorting. Holding $w_i$ fixed in that counterfactual would require a separate equilibrium argument.
\section{Unresolved components}
The sharp results apply to stable quasi-linear preferences, validated offer sets, finite valuation support, and an individual panel or correctly linked conditional experiment. The full research question allows nonlinear utility, many amenities, unemployment, preference shocks, constrained search, and endogenous matching. Those features are not estimated here. In particular, lifetime job security may require an intertemporal expected-utility model rather than the linear amenity index in (1).

The offer-threshold argument and finite-box optimization are standard tools, proved here to expose exactly which inequality claims the data can support. There is no claim that the sign of wage-equivalent inequality has been resolved for any real population.
\begin{thebibliography}{9}
\bibitem{mas} A. Mas (2025). \emph{Non-Wage Amenities}.
RFBerlin discussion paper, sections on identification and the conclusion.
The review motivates caution about wage regressions and job attributes; the restricted box characterization here is self-contained.
\url{https://www.rfberlin.com/wp-content/uploads/2025/08/25051.pdf}.
\end{thebibliography}

\end{document}
